\section{2050 Climate Projections}
\label{sec:projections}

\subsection{NOAA Intermediate SLR Scenario (0.5m Global, 2050)}

Under NOAA's Intermediate scenario, global mean sea level rises 0.5m by 2050 relative
to a 2000 baseline \citep{NOAA_SLR2022}. Gulf Coast regional adjustments for land
subsidence and ocean dynamics produce effective local SLR of 0.6--0.7m on the Louisiana
coast and 0.5--0.55m on the Texas coast. Combined with the independent Louisiana
subsidence rate (1.5 cm/yr over 25 years = 0.375m additional effective lowering of the
land surface), the net effective exposure change for Gulf Coast Louisiana I-10 by 2050
is 0.875--0.975m relative to the 2024 road surface elevation.

We apply this effective SLR to the National Elevation Dataset (1/3 arc-second, 2023
vintage) to compute 2050 SFHA overlay, rerunning the corridor intersection methodology
from Section 3. The 2050 SFHA boundaries for 100-year flood events are expanded to
incorporate the additional 0.875m water column added to current flood-stage elevations.

\subsection{I-10 Gulf Coast Louisiana: 2024 vs. 2050}

At 0.875m effective SLR + subsidence, the 1\% annual exceedance water surface elevation
along the Louisiana reach of I-10 increases by 0.875m. Applying this to the NED elevation
raster: the consecutive SFHA corridor increases from 127 miles to 189 miles, as segments
that are currently marginally above the 100-year flood line (by 0.1--0.5m) fall below
the projected 2050 flood stage. Total SFHA miles increase from 189 to 241. Applying
the D1 formula with 2050 values: $f_c(189) = 10$ (above 60-mi anchor; extended calibration
curve gives 9.8 at 189mi), $f_t(241) = 10$ (above 150-mi anchor; gives 9.9). D1 =
$0.7 \times 9.8 + 0.3 \times 9.9 = 6.86 + 2.97 = 9.1$ (projected).

This represents an increase of 0.7 D1 points (8.4 $\to$ 9.1) driven entirely by land
subsidence and sea level rise. At a 9.1 score, Gulf Coast I-10 Louisiana is not merely
the highest-exposure flood corridor --- it approaches the effective maximum of the D1
scale. No hardening intervention that does not address the structural elevation deficit
can reduce the 2050 score materially.

\subsection{I-10 Gulf Coast Texas: 2024 vs. 2050}

Texas Gulf Coast effective SLR (NOAA regional adjustment + Houston subsidence at 0.8
cm/yr over 25 years = 0.2m additional): net effective exposure increase of 0.7m.
Consecutive SFHA miles: 89 $\to$ 143 miles. Total SFHA miles: 147 $\to$ 198 miles.
Projected D1 = $0.7 \times f_c(143) + 0.3 \times f_t(198)$. Applying the extended
calibration: $f_c(143) = 9.5$, $f_t(198) = 9.9$. D1 = $0.7 \times 9.5 + 0.3 \times
9.9 = 6.65 + 2.97 = 8.8$ (projected, from 7.8 in 2024).

\subsection{I-95 Miami: 2024 vs. 2050}

Miami Beach and Brickell/downtown Miami have among the highest rates of ``sunny day
flooding'' (nuisance flooding at high tide) in the United States, reflecting the
limestone karst geology that allows sea water to infiltrate upward through the substrate.
NOAA's Miami regional SLR at Intermediate scenario: 0.53m by 2050. Combined with coastal
aquifer response, effective 2050 flood stage increase along Miami I-95: 0.6m. Consecutive
SFHA miles: 22 $\to$ 47 miles (the Biscayne Bay approaches and downtown Miami sections
that are currently above the 100-year line fall below it). Total SFHA miles: 61 $\to$
99 miles. Projected D1 = $0.7 \times f_c(47) + 0.3 \times f_t(99)$.
$f_c(47) = 5 + 5 \times (47-15)/(60-15) = 5 + 3.56 = 8.56$; $f_t(99) = 5 + 5 \times
(99-40)/(150-40) = 5 + 2.68 = 7.68$. D1 = $0.7 \times 8.56 + 0.3 \times 7.68 = 5.99 +
2.30 = 7.5$ (projected, from 6.2 in 2024; +1.3 point increase).

\subsection{Priority Order Shift: 2024 vs. 2050}

Table~\ref{tab:2050-projections} compares 2024 and 2050 D1 scores for the top-10
corridors under NOAA Intermediate SLR.

\begin{table}[ht]
\centering
\caption{D1 climate exposure scores: 2024 (current) vs.\ 2050 projected under NOAA
Intermediate SLR (0.5m global, 2050). Gulf Coast corridors incorporate land subsidence
in the effective SLR figure. Winter corridors (Donner, Snoqualmie, Oklahoma) reflect
storm intensity projections rather than SLR.}
\label{tab:2050-projections}
\begin{tabular}{@{}llccc@{}}
\toprule
Rank (2050) & Corridor & D1 (2024) & D1 (2050) & Change \\
\midrule
1 & I-10 Gulf Coast LA  & 8.4 & 9.1 & $+0.7$ \\
2 & I-10 Gulf Coast TX  & 7.8 & 8.8 & $+1.0$ \\
3 & I-80 Donner Pass    & 7.8 & 8.0 & $+0.2$ \\
4 & I-35 Oklahoma       & 7.2 & 7.4 & $+0.2$ \\
5 & I-95 Florida (Miami)& 6.2 & 7.5 & $+1.3$ \\
6 & I-90 Snoqualmie     & 7.1 & 7.3 & $+0.2$ \\
7 & I-5 Siskiyou        & 6.8 & 7.2 & $+0.4$ \\
8 & I-10 Gulf Coast MS  & 6.5 & 7.0 & $+0.5$ \\
9 & I-95 Baltimore      & 4.1 & 5.8 & $+1.7$ \\
10 & I-90 Homestake MT  & 6.4 & 6.5 & $+0.1$ \\
\bottomrule
\end{tabular}
\end{table}

The 2024 priority order places Donner Pass (7.8) in a tie with Gulf Coast Texas (7.8),
slightly below Gulf Coast Louisiana (8.4). By 2050 under Intermediate SLR, Gulf Coast
Louisiana (9.1) and Gulf Coast Texas (8.8) pull away from Donner (8.0) and Snoqualmie
(7.3). I-95 Miami shows the largest proportional increase (+1.3 points), rising from
rank 6 to rank 5.

The most striking entry is I-95 Baltimore (+1.7 points, rising from rank 11 to rank 9):
the Chesapeake Bay coastal sections of I-95 near Baltimore, currently above the 100-year
flood line by a small margin, fall below it by 2050 under NOAA Intermediate SLR, producing
a step-change in SFHA overlay.

\subsection{Infrastructure Lead Time and Authorization Urgency}

Major corridor hardening --- roadway elevation, drainage reconstruction, flood barrier
installation --- requires 8--12 years from congressional authorization to project completion,
based on FHWA project delivery data for major highway reconstruction projects
\citep{IIJA2021}. A project authorized in 2038 reaches completion in 2046--2050. To deliver
Gulf Coast I-10 hardening before the 2050 exposure threshold is reached, authorization
must occur by approximately 2038 --- twelve years from the time of this analysis.

This lead time framing is not rhetorical. The current PROTECT program ($8.7B over five
years) is insufficient in scale to fund Gulf Coast I-10 hardening (estimated \$18--25B
for full-corridor elevation; see D.2 for cost detail). A separate authorization, at the
scale of a Katrina-class infrastructure response, is required if the 2050 D1 = 9.1
exposure is to be addressed before it materializes.
